% ============================================================
\section{Category Partition}
\label{sec:cp}
% ============================================================

\textit{Category Partition} è applicato come tecnica di progettazione
black-box: le classi di equivalenza derivano da documentazione pubblica
e conoscenza di dominio.

\subsection{BrokerImpl}

\subsubsection{\texttt{find(Object oid, boolean validate, FindCallbacks call)}}

\textit{Categorie e partizioni:} oltre a \texttt{oid}, \texttt{validate}, \texttt{call}, sono
modellati come parametri aggiuntivi la presenza dell'entità in cache e nello store — riferite
all'identità effettivamente usata da \texttt{find} in quel momento (l'\texttt{oid} originale, o
quello restituito da \texttt{processArgument} se sostituito) — e i valori restituiti da
\texttt{call.processArgument}/\texttt{call.processReturn}, che pur essendo dichiarati
sull'interfaccia \texttt{FindCallbacks} sembrano modificare il comportamento di \texttt{find}.
Con \texttt{oid} valido si intende un oggetto conforme al tipo usato come identificatore della
classe persistente a cui si riferisce. Classi in Tabella~\ref{tab:find-classi} (Appendice).

\textit{Note di modellazione:} \{istanza non valida\} di \texttt{call} è stata scartata per
assenza di una chiara definizione.
Cache e store sono modellati come due parametri indipendenti, entrambi dipendenti
dall'esistenza di un'identità effettiva: l'\texttt{oid} originale quando \texttt{call} è null,
oppure l'id restituito da \texttt{processArgument} quando \texttt{call} non è null e lo
sostituisce con successo — indipendentemente da quale fosse l'\texttt{oid} di partenza. La
doppia cache \texttt{\_cache}/\texttt{\_transCache} è trattata come un'unica classe, nessuna
fonte ne documenta un effetto diverso su \texttt{find}.
L'id restituito da \texttt{processArgument} (\texttt{@return}: ``the id to look up'') è
dichiarato come l'id da cercare, ma nessuna fonte descrive il meccanismo con cui \texttt{find}
lo utilizza: si assume (ipotesi) che sostituisca l'\texttt{oid} originale come identità su cui
\texttt{find} opera — cache/store si applicano quindi alla nuova identità, non più a quella di
partenza, indipendentemente da come fosse classificato l'\texttt{oid} originale (null, non
valido o valido). Quando invece \texttt{processArgument} ignora l'argomento (ritorna
\texttt{null}, un'istruzione dichiarata distinta da ``il nuovo id è null'') o lancia eccezione,
non esiste alcuna identità effettiva: cache/store collassano anche se l'\texttt{oid} originale
era valido.

\textit{Oracolo:} il coordinamento tra \texttt{processArgument} e \texttt{processReturn} non è
documentato, ma si distinguono due casi. Quando \texttt{processArgument} lancia eccezione,
l'esecuzione si interrompe per semantica standard di Java (eccezione non gestita): si assume che
\texttt{processReturn} non venga raggiunto — le combinazioni restano comunque distinte per
\texttt{processReturn} variabile, per osservare con i test se l'assunzione regge. Quando invece
\texttt{processArgument} ritorna normalmente — sostituendo l'id, o ignorando l'argomento con un
ritorno \texttt{null} — non c'è alcuna interruzione di flusso: si assume che
\texttt{processReturn} venga comunque invocato e che il suo valore determini l'esito di
\texttt{find} (ipotesi basata sul solo \texttt{@return} ``the object to return''), anche quando
l'argomento è stato ignorato. Un ritorno \texttt{null} di \texttt{processReturn} con
\texttt{validate=false} confligge con la garanzia ``mai null'': si assume che quest'ultima
prevalga, essendo dichiarata come proprietà centrale del metodo. Per \texttt{validate=false}
l'ordine cache/hollow/eccezione non è dichiarato esplicitamente: si legge come sequenza di
fallback (cache ha priorità); il successo della creazione hollow non ha alcun appiglio testuale
ed è assunto come esito di default (``hollow'' è comunque terminologia reale del progetto, non
solo prosa del Javadoc: \texttt{PCState.HOLLOW}, ``Hollow; exists in data store''). Quando vale
l'ipotesi che \texttt{processReturn} governi
l'esito (id sostituito con successo, o argomento ignorato), l'oracolo dipende solo da
\texttt{validate} e \texttt{processReturn}, non più da cache/store/hollow — anche quando
l'\texttt{oid} originale era null o non valido.
Le tuple prive di oracolo diretto sono classificate
\emph{inferenza} (esito ricavabile da un testo) o \emph{ipotesi} (nessun appiglio). Su 84 combinazioni valide: 6
dirette, 61 per inferenza, 17 per ipotesi.

\textit{Combinazioni:} $3 (\texttt{oid}) \times 2 \times 2 \mbox{ (cache/store)} \times 2
(\texttt{validate}) \times 2 (\texttt{call}) \times 3 \times 2
(\texttt{processArgument}/\texttt{processReturn}) = 288$.

\textit{Eliminazioni:} solo strutturali (nessuna identità effettiva su cui verificare
cache/store). Quando \texttt{call} è null, l'identità effettiva è l'\texttt{oid} originale:
cache/store collassano ($4 \to 1$) se \texttt{oid} è null o non valido. Quando \texttt{call} non
è null, l'identità effettiva dipende invece solo da \texttt{processArgument}, indipendentemente
da quale fosse l'\texttt{oid} di partenza: se sostituisce con successo, cache/store si applicano
alla nuova identità su tutti e tre i valori di \texttt{oid} (non solo quello valido); se ignora
l'argomento o lancia eccezione, non c'è alcuna identità effettiva e cache/store collassano anche
per \texttt{oid} originariamente valido. Nessun'altra combinazione viene eliminata: le ipotesi
sull'oracolo non giustificano un'esclusione, non c'è motivo a priori di escludere una
combinazione ammissibile perché potrebbe rivelarsi un caso di test significativo. Combinazioni
valide: $84$ (204 eliminate).

\textit{Boundary Value Analysis:} non applicabile, nessuno dei sette parametri ha un range
numerico o una soglia.

\subsubsection{\texttt{lock(Object pc, int level, int timeout, OpCallbacks call)}}

\textit{Categorie e partizioni:} \texttt{pc}: \{null\}, \{non gestito\}, \{gestito, non ancora
persistito\}, \{gestito, presente nello store\} — definite tramite \texttt{PCState}
(\texttt{TRANSIENT}, \texttt{PNEW}, \texttt{PCLEAN}/\texttt{PDIRTY}/\texttt{HOLLOW}
rispettivamente), esposto da \texttt{getPCState()} su \texttt{OpenJPAStateManager} (collegato
strutturalmente a \texttt{call}, terzo parametro di \texttt{OpCallbacks.processArgument}).
\texttt{level}: \{valore noto\}, \{valore non noto\} — \texttt{Broker} estende direttamente
\texttt{LockLevels}, che dichiara tre costanti con Javadoc proprio (\texttt{LOCK\_NONE}=0, ``No
lock''; \texttt{LOCK\_READ}=10, ``Generic read lock level''; \texttt{LOCK\_WRITE}=20, ``Generic
write lock level''). Il primo gruppo copre queste 3 costanti, il secondo qualunque altro valore
\texttt{int}.
\texttt{timeout}: \{altro negativo ($<-1$)\}, \{\texttt{-1}\}, \{0\},
\{positivo\}. \texttt{call}: \{null\}, \{non null\}. \texttt{call.processArgument}* (richiede
\texttt{call} non null): \{ritorna un codice azione (\texttt{ACT\_NONE}/\texttt{ACT\_CASCADE}/
\texttt{ACT\_RUN}, tutte equivalenti)\}, \{lancia \texttt{OpenJPAException}\}. \texttt{OpCallbacks}
ha un solo metodo di callback (nessun equivalente di \texttt{processReturn}). Classi in
Tabella~\ref{tab:lock-classi} (Appendice).

\textit{Note di modellazione:} \texttt{level} è passato per più revisioni. La prima proposta,
\{\texttt{LOCK\_NONE}\}/\{positivo\}, univa in una sola classe ``positivo'' qualunque valore
positivo (costante dichiarata o arbitrario) — stesso problema di equivalenza indebita già
corretto per \texttt{pc}. Una seconda proposta correggeva questo separando \{\texttt{LOCK\_NONE}\},
\{positivo noto\}, \{positivo non noto\}, motivando la separazione di \texttt{LOCK\_NONE} dal
resto solo con la sua diversità semantica (``nessun lock''), lo stesso tipo di criterio già
scartato per \texttt{LOCK\_READ} vs \texttt{LOCK\_WRITE} (semanticamente distinti — lettura vs
scrittura — ma uniti perché nessuna fonte ne documenta un effetto diverso su \texttt{lock}).
Applicando lo stesso criterio in modo coerente, \texttt{LOCK\_NONE}/\texttt{LOCK\_READ}/
\texttt{LOCK\_WRITE} uniscono in un'unica classe \{valore noto\}: \texttt{LockLevels} le dichiara
come tre costanti sorelle, nessuna delle tre documentata con un effetto diverso sull'esito di
\texttt{lock} — la partizione corretta è quindi a due classi (\{valore noto\}/\{valore non
noto\}), non a tre.

Una revisione successiva ha inoltre rimosso un riferimento, presente in una versione precedente
di questa analisi, a \texttt{LockManager.lock} e a \texttt{MixedLockLevels} come fonti per il
dominio di \texttt{level}. \texttt{LockManager} è raggiungibile da \texttt{Broker} (tramite
\texttt{StoreContext.getLockManager()}, la stessa distanza strutturale già usata per
\texttt{StoreManager}), ma il suo metodo \texttt{lock(OpenJPAStateManager, int, int, Object)} è
un metodo diverso, su un'interfaccia diversa, senza alcuna relazione di delega dichiarata con
\texttt{Broker.lock(Object, int, int, OpCallbacks)}: usare il suo Javadoc (``one of the lock
constants defined in \texttt{LockLevels}, or a custom level'') per definire il dominio di
\texttt{Broker.lock} avrebbe richiesto assumere che quest'ultimo deleghi internamente al primo,
cosa verificabile solo leggendo l'implementazione (violazione BB). \texttt{MixedLockLevels} era
una fonte ancora più debole: non è estesa né da \texttt{Broker} né da \texttt{LockManager} (che
estende solo \texttt{LockLevels}), e le sue costanti compaiono nel codice sorgente solo dentro
\texttt{FetchConfigurationImpl}, una classe di implementazione concreta priva di qualunque
connessione strutturale con \texttt{BrokerImpl} — una citazione non ammissibile, rimossa.
L'unica fonte strutturalmente diretta per i valori noti di \texttt{level} resta quindi
\texttt{LockLevels}, estesa direttamente da \texttt{Broker}: 3 costanti, non 8.

Va inoltre precisato che un valore non noto (es. \texttt{level=10000}) non ha un significato di
dominio interpretabile da chi progetta i test: \texttt{level} è dichiarato semplicemente
\texttt{int}, senza altra restrizione nel Javadoc di \texttt{Broker.lock}, e nessuna fonte
strutturalmente connessa a \texttt{BrokerImpl} descrive un significato per valori diversi dalle 3
costanti di \texttt{LockLevels} (non esiste, ad esempio, una scala documentata di livelli
intermedi). La classe \{valore non noto\} è quindi trattata a scatola nera come ``qualunque
intero non tra le costanti dichiarate''. In assenza di qualunque fonte, si assume (\emph{ipotesi})
un comportamento diverso da quello dei livelli noti, verosimilmente un'eccezione, anziché
identico. Le
classi di \texttt{pc} sostituiscono un'iniziale coppia \{non valido\}/\{valido\} basata su
conoscenza di dominio generica, poco verificabile in pratica.
\texttt{PCState} documenta esplicitamente il significato di ciascuno stato (\texttt{TRANSIENT}:
``Transient; unmanaged instance''; \texttt{PNEW}: ``Persistent-New''; \texttt{PCLEAN}/
\texttt{PDIRTY}: ``Persistent-Clean''/``Persistent-Dirty''), verificabile in un test mockando
\texttt{getPCState()}. Il livello di lock corrente di \texttt{pc} non è tracciato come parametro
a sé: la postcondizione dichiarata (``ensure locked at the given level'') non dipende dallo
stato di partenza. \texttt{PCState} dichiara in totale 23
costanti; oltre alle 4 già usate, solo \texttt{HOLLOW} (``Hollow; exists in data store'') ha un
appiglio testuale diretto rilevante qui, e viene incluso nella classe \{gestito, presente nello
store\} per lo stesso criterio già usato per \texttt{PCLEAN}/\texttt{PDIRTY} — confermato anche
da \texttt{StoreManager.load} (\texttt{StoreManager} raggiungibile via
\texttt{StoreContext.getStoreManager()}), che distingue ``a hollow state manager''
(``fields must be loaded'') da ``a non-hollow state manager''. Le restanti 18
(varianti non-transazionali, \texttt{TCLEAN}/\texttt{TDIRTY}/\texttt{TLOADED}, stati
\texttt{Embedded}, combinazioni fini di flush/delete) non sono modellate come classi separate:
le non-transazionali richiederebbero una dimensione SUT aggiuntiva non giustificata da alcuna
fonte specifica su \texttt{lock}; le transient-qualificate (\texttt{TCLEAN}/\texttt{TDIRTY}/
\texttt{TLOADED}) condividono il nome con \texttt{TRANSIENT} ma non il suo Javadoc: solo
\texttt{TRANSIENT} dichiara esplicitamente ``unmanaged instance'', le altre tre hanno un
Javadoc di una riga (``Transient-Clean/-Dirty/-Loaded'') privo di quella dicitura — condividere
una parola nel nome non basta ad assumere lo stesso significato; per confermarlo servirebbe
leggere le classi di transizione di stato (WB), quindi restano fuori, non incluse in
\{non gestito\}; gli stati \texttt{Embedded} riguardano value object, non risulta
che \texttt{lock} li accetti come \texttt{pc}; le combinazioni di flush/delete sono casi limite
della pipeline di cancellazione, senza alcuna fonte che li distingua per \texttt{lock}. Un
eventuale esito di timeout scaduto sotto contesa concorrente non è documentato da
\texttt{Broker.lock} stesso e non è incluso nella combinatoria: richiede comunque una contesa
concorrente reale (un secondo thread/broker che mantiene il lock su \texttt{pc} oltre il
\texttt{timeout}), quindi un test di integrazione con più componenti reali, non un unit test con
collaboratori mockati — candidato per una fase di test di integrazione separata (dimensione
Level: Unit/Integration).

\textit{Ipotesi sull'oracolo:} per \texttt{pc} null o \texttt{TRANSIENT} (non gestito), quando
\texttt{processArgument} non lancia essa stessa un'eccezione, nessuna fonte documenta l'esito: si
assume un'eccezione (tipo non specificato), \emph{ipotesi} senza appiglio testuale. Per
\texttt{pc} gestito (\texttt{PNEW}, oppure \texttt{PCLEAN}/\texttt{PDIRTY}/\texttt{HOLLOW}) con
\texttt{level} \emph{non noto}, si assume anch'esso un'eccezione (tipo non specificato),
\emph{ipotesi} (stesso giudizio di progettazione motivato sopra). Resta invece \texttt{pc}
gestito con \texttt{level}
\emph{noto}: per \texttt{PCLEAN}/\texttt{PDIRTY}/\texttt{HOLLOW} (gestito, presente nello store)
la postcondizione (\texttt{getLockLevel(pc) == level}, verificata con l'inspector
\texttt{StoreContext.getLockLevel}) è \emph{diretta}. Per \texttt{PNEW} (gestito, non ancora
persistito) non è documentato se ``ensure'' valga anche prima del flush verso lo store: si
assume, per analogia, che la stessa postcondizione valga comunque — \emph{ipotesi}, non
\emph{diretta}, perché estende un caso documentato a uno strutturalmente diverso. La premessa
``non ancora persistito'' non è solo dedotta dal nome \texttt{PNEW}: \texttt{StoreManager.flush}
(\texttt{Collection<OpenJPAStateManager> sms}) distingue esplicitamente stati che ``do not
require data store action, such as persistent-clean instances'' da ``New instances without an
assigned object id'', che lo ricevono proprio durante il flush — quindi nessuna riga reale da
bloccare finché l'istanza non è persistita. Non si modella un'ulteriore sotto-classe ``\texttt{PNEW}
già flushato'': il metodo \texttt{PCState.flush(StateManagerImpl context)} (dichiarato sulla
classe base, Javadoc ``return the proper lifecycle state on flush. Returns the \texttt{this}
pointer by default'') conferma che il flush può cambiare lo stato restituito da
\texttt{getPCState()} — quale sia lo stato risultante specifico per \texttt{PNEW} richiederebbe
leggere l'override in \texttt{PNewState} (WB, stessa esclusione già applicata a
\texttt{TCLEAN}/\texttt{TDIRTY}/\texttt{TLOADED}). Non serve saperlo: un'istanza osservata come
\texttt{PNEW} è per costruzione ``non ancora flushata'' in quel momento, altrimenti
\texttt{getPCState()} restituirebbe già un altro valore — verosimilmente uno già coperto dalla
classe \{gestito, presente nello store\}. Quando \texttt{call.processArgument} lancia
eccezione, si assume la propagazione da \texttt{lock} (\emph{inferenza}: propagazione delle
eccezioni non gestite, stessa proprietà usata per \texttt{find}) — indipendentemente dalla
classe di \texttt{pc} o di \texttt{level}: è un fatto di semantica Java, non legato allo stato
dell'istanza, stesso trattamento già usato per \texttt{detachAll}/\texttt{attachAll}/
\texttt{persist}. Quando \texttt{pc} è gestito con \texttt{level} \emph{non noto} e
\texttt{processArgument} ritorna un codice azione (senza lanciare eccezione), il valore atteso
resta lo stesso dell'intera riga (eccezione, per il giudizio di progettazione appena descritto):
il codice azione non cambia l'ipotesi, dato che a monte non ci si aspetta comunque che il
livello venga applicato. Quando invece \texttt{level} è \emph{noto} e \texttt{processArgument}
ritorna un codice azione, l'effetto su \texttt{lock} non è
documentato da \texttt{OpCallbacks} (interfaccia generica, condivisa da tutte le operazioni,
senza semantica specifica per \texttt{lock}): il valore atteso resta identico a quello della
riga gemella con \texttt{call} null (stesso \texttt{pc}/\texttt{level}/\texttt{timeout}) — si
assume che ``ensure locked at the given level'' resti comunque garantito, senza essere alterato
dal codice azione — ma il tier scende comunque a \emph{ipotesi} (o resta \emph{ipotesi}, se la
gemella già lo era) perché quel codice azione introduce un'interazione non documentata che
potrebbe, in teoria, cambiare l'esito. Il \texttt{timeout=-1} è documentato esattamente da
\texttt{Broker.lock} (``-1 for no limit''); nessuna fonte strutturalmente connessa a
\texttt{Broker.lock} — ci si affida solo alla documentazione di \texttt{lock()} stesso, non a
\texttt{LockManager} — documenta il comportamento per un negativo diverso da \texttt{-1}. Per
\texttt{pc} gestito con \texttt{level} \emph{noto}, la riga con \texttt{timeout} esattamente
\texttt{-1} resta \emph{diretta} (o \emph{ipotesi} se già scesa per altri motivi, es.
\texttt{PNEW} o azione non eccezionale); per la classe \{altro negativo ($<-1$)\} si assume
invece un'eccezione (tipo non specificato), \emph{ipotesi}: dato che solo il valore esatto
\texttt{-1} è documentato, un negativo diverso è più plausibilmente respinto che trattato allo
stesso modo. Per le righe dove il tier era già
\emph{ipotesi} per un motivo indipendente da \texttt{timeout} (\texttt{pc} non
gestito/\texttt{null}, o \texttt{level} non noto) la classe \{altro negativo\} non cambia
nulla — l'oracolo era già ``eccezione''; per le righe con propagazione di eccezione da
\texttt{processArgument} resta \emph{inferenza} indipendentemente da \texttt{timeout} (fatto di
semantica Java, non legato al valore numerico).

\textit{Correzione di una partizione non esaustiva su \texttt{timeout}:} una revisione
successiva del parametro \texttt{timeout} ha rilevato che la partizione iniziale
$\{\texttt{-1}\}, \{0\}, \{\texttt{positivo}\}$ non copriva l'intero dominio dei valori
ammissibili per un \texttt{int}: ogni valore negativo diverso da \texttt{-1} (es. \texttt{-2},
\texttt{-100}) non apparteneva a nessuna delle tre classi. L'errore nasceva dal confondere due
fasi distinte del metodo Category Partition: la definizione delle classi di equivalenza
(partizionare il dominio di un parametro, verificando che l'unione delle classi copra l'intero
dominio) e il Boundary Value Analysis (scegliere, \emph{dentro} una classe già definita, un
valore rappresentativo al confine). Trattare \texttt{-1} come l'unico rappresentante scelto per
un'ipotetica classe ``negativo'' avrebbe fatto collassare silenziosamente queste due fasi: la
classe ``negativo'' non era mai stata definita esplicitamente al di là del singolo valore
\texttt{-1}, quindi tutti gli altri negativi restavano semplicemente non testati, senza che una
decisione di eliminazione lo giustificasse esplicitamente. La correzione consiste
nell'aggiungere la classe mancante \{altro negativo ($<-1$)\}, completando la partizione (ora
esaustiva su tutto il dominio \texttt{int}); il valore rappresentativo scelto per questa classe
via BVA è \texttt{-2}, il confine adiacente a \texttt{-1} (l'altro estremo della classe è
illimitato verso $-\infty$ e non ha un confine documentato da esplorare). La lezione generale,
applicata poi a ritroso su tutti gli altri parametri numerici già modellati per gli altri 6
metodi (nessun altro caso analogo trovato): per ogni parametro, prima verificare che le classi
di equivalenza già definite coprano esaustivamente l'intero dominio dichiarato, e solo dopo, per
ciascuna classe, applicare il BVA per scegliere un valore di confine — mai lasciare che un
singolo valore scelto per comodità (perché documentato, o perché ``ovvio'') faccia da
sostituto implicito per un'intera classe più ampia mai definita esplicitamente.

\textit{Correzione di una classe di equivalenza non giustificata su
\texttt{call.processArgument}:} una revisione ha rilevato che \{\texttt{ACT\_NONE}\} e
\{azione: \texttt{ACT\_CASCADE}/\texttt{ACT\_RUN}\} erano tenute come due classi separate senza
alcuna fonte che le distinguesse: \texttt{OpCallbacks.java} non ha Javadoc per nessuna delle tre
costanti (verificato leggendo il file), esattamente come nessuna fonte distingueva
\texttt{LOCK\_NONE} da \texttt{LOCK\_READ}/\texttt{LOCK\_WRITE} per \texttt{level} — stesso
errore, stessa correzione: le tre costanti uniscono in un'unica classe \{ritorna un codice
azione\}. La partizione di \texttt{call.processArgument} passa quindi da tre classi a due:
\{ritorna un codice azione\}, \{lancia \texttt{OpenJPAException}\}.

\textit{Combinazioni:} $4 (\texttt{pc}) \times 2 (\texttt{level}) \times 4 (\texttt{timeout})
\times 2 (\texttt{call}) \times 2 (\texttt{processArgument}) = 128$.

\textit{Eliminazioni:} solo strutturale, quando \texttt{call} è null non esiste
\texttt{processArgument} da invocare (collassa da 2 a 1). Combinazioni valide:
$4 \times 2 \times 4 \times (1+2) = 96$ (32 eliminate).

\textit{Boundary Value Analysis:} \texttt{timeout} ha un confine documentato esplicitamente a
\texttt{-1} (fonte primaria \texttt{Broker.lock}): la classe \{\texttt{-1}\} è un valore
sentinella singolo, non un range, quindi coincide col proprio stesso rappresentante. La classe
\{altro negativo ($<-1$)\} ha un confine noto solo sul lato adiacente a \texttt{-1}
(rappresentante BVA: \texttt{-2}) e nessun confine documentato verso $-\infty$. Le classi
\{0\} e \{positivo\} condividono il confine $0$/$1$, già coperto scegliendo \texttt{0} e un
valore positivo qualunque come rappresentanti delle rispettive classi. \texttt{level} non ha
un confine numerico documentato con effetto diverso su \texttt{lock} (nessun limite superiore,
nessuna soglia dichiarata) — la partizione \{noto\}/\{non noto\} è una classe di equivalenza
Category Partition (da distinzione testuale), non un confine BVA: non si applica BVA su un
dominio senza soglia documentata.

\textit{Oracolo:} su 96 combinazioni valide, 3 hanno oracolo diretto (\texttt{pc}
\texttt{PCLEAN}/\texttt{PDIRTY}/\texttt{HOLLOW}, \texttt{call} null, \texttt{level} \emph{noto},
\texttt{timeout} esattamente \{\texttt{-1}\}, \{0\} o \{positivo\} — non \{altro negativo\}), 32
per inferenza (propagazione eccezione di \texttt{processArgument} su tutte le combinazioni di
\texttt{pc}/\texttt{level}/\texttt{timeout}: la propagazione è un fatto strutturale di Java,
indipendente da entrambi), 61 per ipotesi: 52 con oracolo ``eccezione'' (\texttt{pc} non gestito
o \texttt{null}; \texttt{pc} gestito con \texttt{level} \emph{non noto}; oppure \texttt{pc}
gestito con \texttt{level} \emph{noto} e \texttt{timeout} \{altro negativo\} — nei due casi
``non noto''/``altro negativo'' un giudizio di progettazione, non un fatto documentato) e 9 con
oracolo ``postcondizione'' (\texttt{pc} \texttt{PNEW} con \texttt{level} \emph{noto} e
\texttt{timeout} $\neq$ \{altro negativo\}, o \texttt{call} con azione non eccezionale su
\texttt{pc} gestito con \texttt{level} \emph{noto} e \texttt{timeout} $\neq$ \{altro negativo\}).
Materializzate per esteso in Appendice.

\subsubsection{\texttt{detachAll(Collection objs, OpCallbacks call)}}

\textit{Categorie e partizioni:} \texttt{objs}: \{null\}, \{vuota\}, \{un elemento gestito\},
\{un elemento non gestito\}, \{elementi misti (un gestito e un non gestito insieme)\} —
definite tramite \texttt{PCState} (stesso costrutto usato per \texttt{lock}/\texttt{persist}):
\{non gestito\} è \texttt{TRANSIENT} (``unmanaged instance''), \{gestito\} è un rappresentante
per qualunque altro stato, non distinto ulteriormente perché \texttt{detachAll} non documenta un
effetto diverso tra sotto-stati gestiti (rappresentante per la cardinalità ``non vuota'', per
coerenza con la logica CP di classi di equivalenza piuttosto che dimensioni esaustive). La
classe ``elementi misti'' testa l'atomicità non documentata dell'operazione su collezione:
nessuna fonte dice se un elemento non gestito faccia fallire l'intera chiamata o solo il proprio
elemento, lasciando intatti gli altri. \texttt{call}: \{null\}, \{non null\}. \texttt{call.processArgument}*
(richiede \texttt{call} non null): \{ritorna un codice azione\}, \{lancia
\texttt{OpenJPAException}\} (stesso tipo \texttt{OpCallbacks} di \texttt{lock}, stessa unione
di \texttt{ACT\_NONE}/\texttt{ACT\_CASCADE}/\texttt{ACT\_RUN} in un'unica classe: nessuna delle
tre ha Javadoc, nessuna fonte le distingue). Classi in Tabella~\ref{tab:detachall-classi}
(Appendice).

\textit{Note di modellazione:} le classi di \texttt{objs} sostituiscono un'iniziale coppia
\{un elemento non valido\}/\{un elemento valido\}, la cui etichetta non coincideva col criterio
usato per l'oracolo (``gestito''/``non gestito''): stesso problema di deriva terminologica già
corretto per \texttt{pc} in \texttt{lock}/\texttt{persist}, qui risolto allineando l'etichetta al
criterio, senza bisogno di un nuovo costrutto (già disponibile tramite \texttt{PCState}). Il
\texttt{@return} di \texttt{detachAll} (``the detached instances'') non definisce da solo cosa
significhi ``detached'': la definizione è nel Javadoc di \texttt{isDetached} (``an object that
can be reattached... via \texttt{Broker\#attach}''), un metodo diverso. L'oracolo per l'elemento
gestito combina le due fonti (oracolo per \emph{inspector}, già usato per \texttt{find}/
\texttt{cache}): per ogni elemento \texttt{e} del risultato si verifica
\texttt{isDetached(e) == true}. Per questo nessuna riga di questo metodo è classificata come
\emph{diretta}: ogni oracolo utile richiede di collegare almeno due fonti. Il flush automatico
descritto per l'overload \texttt{detachAll(OpCallbacks)} non si applica qui: l'overload a 2
parametri scelto per il test non lo menziona nel proprio Javadoc.

\textit{Ipotesi sull'oracolo:} \texttt{objs} null o con un elemento non gestito
(\texttt{TRANSIENT}): nessun appiglio testuale, si assume un'eccezione (tipo non specificato),
\emph{ipotesi}. \texttt{objs} con elementi misti (un gestito e un non gestito insieme): nessuna
fonte documenta l'atomicità dell'operazione; si assume, per analogia col caso di solo elemento
non gestito, che l'intera chiamata fallisca con eccezione — \emph{ipotesi}, estende un caso già
ipotetico a uno strutturalmente diverso (non un'estensione di un caso diretto). \texttt{objs}
vuota: un array vuoto, dedotto per corrispondenza di cardinalità col \texttt{@return} plurale,
\emph{inferenza}. \texttt{processArgument} che lancia eccezione: propagazione da
\texttt{detachAll} (\emph{inferenza}, stessa proprietà generale usata per \texttt{find}/
\texttt{lock}). Per il codice azione: nessuna descrizione dell'effetto su
\texttt{detachAll} in \texttt{OpCallbacks} (interfaccia generica, non specifica); si assume che
il comportamento di base non cambi, \emph{ipotesi} senza un verbo di garanzia assoluto a
supporto (a differenza di \texttt{lock}).

\textit{Combinazioni:} $5 (\texttt{objs}) \times 2 (\texttt{call}) \times 2
(\texttt{processArgument}) = 20$.

\textit{Eliminazioni:} solo strutturale, \texttt{call} null collassa \texttt{processArgument} da
2 a 1. Combinazioni valide: $5 \times (1+2) = 15$ (5 eliminate).

\textit{Boundary Value Analysis:} non applicabile, nessun parametro numerico con soglia
documentata.

\textit{Oracolo:} su 15 combinazioni valide, 7 per inferenza (cardinalità/inspector,
propagazione eccezione), 8 per ipotesi (\texttt{objs} inammissibile o mista, o \texttt{call}
con codice azione). Nessuna riga diretta. Materializzate per esteso in Appendice.

\subsubsection{\texttt{attachAll(Collection objs, boolean copyNew, OpCallbacks call)}}

\textit{Categorie e partizioni:} \texttt{objs}: \{null\}, \{vuota\}, \{un elemento
precedentemente \texttt{detached}\}, \{un elemento nuovo\}, \{un elemento già gestito da questo
broker\} (caso non coperto dai due rami documentati, si veda sotto), \{un elemento non valido\}
(\texttt{PCState.PDELETED}, ``Persistent-Deleted'' — marcato per cancellazione, stato non
considerato da nessuno dei rami documentati), \{elementi misti (un valido e un non valido
insieme)\}. \texttt{copyNew}: \{true\}, \{false\}. \texttt{call}: \{null\}, \{non null\}.
\texttt{call.processArgument}* (richiede \texttt{call} non null): \{ritorna un codice azione\},
\{lancia \texttt{OpenJPAException}\} (stessa unione già applicata a \texttt{lock}/
\texttt{detachAll}: nessuna delle tre costanti ha Javadoc). Classi in
Tabella~\ref{tab:attachall-classi} (Appendice).

\textit{Note di modellazione:} il Javadoc descrive esattamente due rami per lo stato di un
elemento di \texttt{objs}: precedentemente distaccato (merge) o nuovo (persistito). Un elemento
già gestito dallo stesso broker (né distaccato né nuovo) non rientra in nessuno dei due — gap
già individuato in fase di analisi documentale, mantenuto come classe a sé perché descrive uno
stato realmente costruibile e distinto, non una sotto-classificazione inventata. Il flag
\texttt{getPostLoadOnMerge}/\texttt{setPostLoadOnMerge} (default \texttt{false}) condiziona un
evento di lifecycle durante il merge, non il valore di ritorno: non tracciato come dimensione,
irrilevante per l'oracolo scelto (il valore restituito, non gli effetti collaterali). A
differenza di \texttt{detachAll}/\texttt{lock}/\texttt{persist}, qui \texttt{TRANSIENT} non
rappresenta il caso non valido: coincide semanticamente col ramo documentato ``nuovo'' (un'
istanza mai gestita, introdotta per la prima volta tramite \texttt{attachAll}). Il caso non
valido è quindi definito su un altro stato non documentato per questo metodo:
\texttt{PDELETED}, un'istanza già marcata per cancellazione — concettualmente incompatibile con
un'operazione di attach/merge, ma nessuna fonte lo conferma esplicitamente. La classe
``elementi misti'' testa, come per \texttt{detachAll}, l'atomicità non documentata
dell'operazione su collezione.

\textit{Ipotesi sull'oracolo:} ramo ``precedentemente \texttt{detached}'' (\emph{diretta}):
l'istanza persistente viene aggiornata con le modifiche, indipendente da \texttt{copyNew} (il
cui \texttt{@param} lo scopa esplicitamente alle sole istanze nuove). Ramo ``nuovo''
(\emph{diretta}): istanza persistita come nuova; \texttt{copyNew=true} implica una copia
dell'originale, \texttt{copyNew=false} no, per dichiarazione esplicita del parametro. Ramo
``già gestito'' (\emph{ipotesi}, nessun appiglio): si assume, per analogia col ramo merge,
nessuna duplicazione. \texttt{objs} non valido (\texttt{PDELETED}): nessuna fonte documenta
l'esito, si assume un'eccezione (\emph{ipotesi}). \texttt{objs} con elementi misti: si assume,
per analogia col caso di solo elemento non valido, che l'intera chiamata fallisca con eccezione
(\emph{ipotesi}, estende un caso già ipotetico). \texttt{objs} null (\emph{ipotesi}) ed
eccezione di \texttt{processArgument} (\emph{inferenza}, propagazione) seguono lo stesso
trattamento già usato per gli altri metodi. Il codice azione: nessuna descrizione
dell'effetto su questo metodo, si assume che il comportamento di base non cambi
(\emph{ipotesi}).

\textit{Combinazioni:} $7 (\texttt{objs}) \times 2 (\texttt{copyNew}) \times 2 (\texttt{call})
\times 2 (\texttt{processArgument}) = 56$.

\textit{Eliminazioni:} solo strutturale, \texttt{call} null collassa \texttt{processArgument} da
2 a 1. Combinazioni valide: $7 \times 2 \times (1+2) = 42$ (14 eliminate).

\textit{Boundary Value Analysis:} non applicabile, nessun parametro numerico con soglia
documentata.

\textit{Oracolo:} su 42 combinazioni valide, 4 dirette (rami merge/nuovo con \texttt{call}
null), 16 per inferenza (cardinalità vuota, propagazione eccezione), 22 per ipotesi
(\texttt{objs} inammissibile, già gestito, non valido, o misto, o \texttt{call} con codice
azione). Materializzate per esteso in Appendice.

\subsubsection{\texttt{newInstance(Class cls)}}

\textit{Categorie e partizioni:} unico parametro formale, nessuno stato SUT aggiuntivo
documentato. \texttt{cls}: \{null\}, \{interfaccia con metodi JavaBeans\}, \{classe astratta con
metodi JavaBeans\}, \{tipo managed concreto\}, \{interfaccia senza metodi JavaBeans\}, \{tipo
non managed, non interfaccia\}. Classi in Tabella~\ref{tab:newinstance-classi} (Appendice).

\textit{Ipotesi sull'oracolo:} i primi tre casi non-null hanno oracolo \emph{diretto}: i due rami
descritti esplicitamente nel Javadoc (creazione di un'implementazione concreta per
interfacce/classi astratte con convenzione JavaBeans; istanza diretta per un tipo managed) più il
\texttt{@throws IllegalArgumentException if cls is not a managed type or interface} per il
sesto caso. Il quinto caso (interfaccia senza convenzione JavaBeans) espone un'ambiguità
testuale reale: il ramo di creazione richiede esplicitamente la convenzione JavaBeans, ma la
clausola \texttt{@throws} esenta genericamente ``interfaccia'' senza ripetere quel vincolo.
Nessun meccanismo di creazione valido si applicherebbe comunque: si assume un'eccezione
(\emph{ipotesi}). \texttt{cls} null non è documentato: si assume un'eccezione, stesso
trattamento degli altri parametri obbligatori (\emph{ipotesi}).

\textit{Combinazioni:} un solo parametro, 6 classi $=$ 6 combinazioni, nessuna eliminazione.

\textit{Boundary Value Analysis:} non applicabile, nessun parametro numerico.

\textit{Oracolo:} 4 dirette, 2 per ipotesi. Materializzate per esteso in Appendice.

\subsubsection{\texttt{isDetached(Object obj, boolean find)}}

\textit{Categorie e partizioni:} \texttt{obj}: \{null\}, \{non gestito\}, \{detached\}, \{non
detached (gestito, persistente)\} — le ultime due sono lo stato reale dell'oggetto, non l'esito
del metodo: quello che il metodo deve determinare correttamente. \{non gestito\} è
\texttt{PCState.TRANSIENT} (``unmanaged instance''), stesso costrutto usato per \texttt{lock}/
\texttt{detachAll}/\texttt{persist} — la distinzione valido/non valido su \texttt{obj} è quindi
realizzata tramite \texttt{PCState}, con l'eccezione di \{detached\}: tra le 23 costanti di
\texttt{PCState} non ne esiste una che rappresenti ``detached''; questo stato è definito solo
dal Javadoc di \texttt{Broker.isDetached(Object)} (citato in
Sezione~\ref{sec:brokerimpl-analisi-doc}), non da \texttt{PCState}.
\texttt{find}: \{true\}, \{false\}. Classi in Tabella~\ref{tab:isdetached-classi} (Appendice).

\textit{Note di modellazione:} il Javadoc di \texttt{find} (``as a last resort this method will
check whether or not the provided object exists in the DB'') descrive un meccanismo di verifica
aggiuntivo attivato solo quando un controllo più economico non è dirimente — condizione interna
non osservabile in BB. Nel modello adottato, \texttt{find} non cambia l'oracolo previsto per
nessuna delle quattro classi di \texttt{obj} (il \texttt{@return} è incondizionato: vero se
detached, falso altrimenti); resta comunque un parametro formale testato su entrambi i valori,
per costruzione, non eliminato.

\textit{Ipotesi sull'oracolo:} \texttt{obj} detached $\to$ \texttt{true}; \texttt{obj} non
detached $\to$ \texttt{false}: entrambe \emph{dirette}, dal \texttt{@return} esplicito (``True if
the provided obj is detached, false otherwise''). \texttt{obj} null o non gestito: non
documentato; essendo un metodo predicato (query, non un'operazione), si assume una risposta
negativa piuttosto che un'eccezione (\emph{ipotesi}).

\textit{Combinazioni:} $4 (\texttt{obj}) \times 2 (\texttt{find}) = 8$, nessuna eliminazione (ogni
parametro ha sempre un referente).

\textit{Boundary Value Analysis:} non applicabile, nessun parametro numerico.

\textit{Oracolo:} 4 dirette, 4 per ipotesi. Materializzate per esteso in Appendice.

\subsubsection{\texttt{persist(Object pc, Object id, OpCallbacks call)}}

\textit{Categorie e partizioni:} \texttt{pc}: \{null\}, \{non gestito\}, \{gestito\} — definite
tramite \texttt{PCState} (\texttt{TRANSIENT} per il primo caso, \texttt{PNEW}/\texttt{PCLEAN}/
\texttt{PDIRTY}/\texttt{HOLLOW} indistinte per il secondo, stesso costrutto già usato per
\texttt{lock}), esposto da \texttt{getPCState()} su \texttt{OpenJPAStateManager}. A differenza di
\texttt{lock}, qui non si distinguono tra loro gli stati già gestiti: nessuna fonte descrive un
comportamento diverso di \texttt{persist} tra \texttt{PNEW}, \texttt{PCLEAN}/\texttt{PDIRTY} e
\texttt{HOLLOW}, tutti accomunati dall'essere già gestiti (nessuno stato manager da assegnare ex
novo). Delle restanti 18 costanti di \texttt{PCState} (varianti non-transazionali, transient-
qualificate, \texttt{Embedded}, casi limite di flush/delete), nessuna ha un appiglio testuale
diverso da quello già usato per \texttt{lock}: stessa esclusione, stesse motivazioni. \texttt{id}: \{null\}, \{non null\} —
dominio esplicitamente dichiarato nel Javadoc del parametro. \texttt{call}: \{null\}, \{non
null\}. \texttt{call.processArgument}* (richiede \texttt{call} non null): \{ritorna un codice
azione\}, \{lancia \texttt{OpenJPAException}\} (stessa unione già applicata a \texttt{lock}/
\texttt{detachAll}/\texttt{attachAll}: nessuna delle tre costanti ha Javadoc). Classi in
Tabella~\ref{tab:persist-classi} (Appendice).

\textit{Note di modellazione:} le classi di \texttt{pc} sostituiscono un'iniziale coppia
\{non valido\}/\{valido\} basata su conoscenza di dominio generica, stesso problema già
riscontrato e corretto per \texttt{lock}. L'overload scelto (a 3 parametri, con \texttt{id}) si
distingue esplicitamente da ``altre operazioni di persist'' per non cascare immediatamente sui
campi \texttt{CASCADE\_IMMEDIATE}: comportamento di cascata sui campi correlati di \texttt{pc},
non sul valore di ritorno testato — non tracciato come dimensione separata, stessa scelta già
applicata alle altre configurazioni di stato collaterali (\texttt{DetachState},
\texttt{PostLoadOnMerge} per i metodi precedenti).

\textit{Ipotesi sull'oracolo:} \texttt{pc} \texttt{TRANSIENT} (non gestito), con \texttt{id} null
o non null: entrambe \emph{dirette}, dal \texttt{@return} (``the state manager for the newly
persistent instance'') combinato col \texttt{@param id} (``may be null for default'') — questo è
il caso descritto dal Javadoc, un'istanza non gestita che diventa persistente. \texttt{pc} già
\texttt{gestito} (\texttt{PNEW}/\texttt{PCLEAN}/\texttt{PDIRTY}/\texttt{HOLLOW}): nessuna fonte
documenta l'esito di \texttt{persist} su un'istanza già persistente; si assume, per analogia, che valga comunque la
stessa postcondizione (\emph{ipotesi}, non \emph{diretta}, per lo stesso motivo già visto in
\texttt{lock} per \texttt{PNEW}: estende un caso documentato a uno strutturalmente diverso).
\texttt{pc} null: non documentato, si assume un'eccezione (\emph{ipotesi}).
\texttt{processArgument} che lancia eccezione: propagazione (\emph{inferenza}, stessa proprietà
generale già usata per gli altri metodi). Il codice azione: nessuna descrizione
dell'effetto su questo metodo, si assume che il comportamento di base non cambi (\emph{ipotesi}).

\textit{Combinazioni:} $3 (\texttt{pc}) \times 2 (\texttt{id}) \times 2 (\texttt{call}) \times 2
(\texttt{processArgument}) = 24$.

\textit{Eliminazioni:} solo strutturale, \texttt{call} null collassa \texttt{processArgument} da
2 a 1. Combinazioni valide: $3 \times 2 \times (1+2) = 18$ (6 eliminate).

\textit{Boundary Value Analysis:} non applicabile, nessun parametro numerico con soglia
documentata.

\textit{Oracolo:} su 18 combinazioni valide, 2 dirette (\texttt{pc} \texttt{TRANSIENT},
\texttt{call} null), 6 per inferenza (propagazione eccezione), 10 per ipotesi (\texttt{pc} null
o già \texttt{gestito}, o \texttt{call} con codice azione). Materializzate per esteso in
Appendice.

\subsection{LoginDialog}

% analisi da svolgere